\chapter{JC Polyomavirus PCR Test}

\section*{What Is This Test?}
JC polyomavirus (JCPyV) PCR detects viral DNA in a clinical specimen. JCPyV commonly remains latent in healthy people, but reactivation can cause progressive multifocal leukoencephalopathy (PML), a serious demyelinating brain infection, or less commonly other disease in severely immunocompromised patients.

\section*{Alternative Names}
JC Virus PCR; JCV DNA PCR; JC Polyoma Virus DNA; Quantitative JCPyV PCR

\section*{How It Is Performed}
Viral DNA is detected by real-time polymerase-chain-reaction testing of cerebrospinal fluid, plasma, urine, or another specimen selected for the clinical question. Quantitative results are usually reported as copies or international units per millilitre when a validated quantitative assay is available.

\section*{Why It Is Ordered}
\begin{itemize}
  \item Evaluation of suspected PML in a patient with compatible neurologic symptoms and immune suppression
  \item Supporting assessment of JCPyV reactivation in selected transplant, hematologic, or immunomodulated patients
  \item Monitoring viral DNA when specialist management requires serial testing
\end{itemize}

\section*{Interpretation and Limitations}
A positive cerebrospinal-fluid result supports JCPyV infection of the central nervous system in the appropriate clinical setting, but a negative result does not exclude PML, particularly early in disease or when the viral load is low. Repeat testing with a sensitive assay and brain MRI may be required; MRI is an imaging study and is not part of this test. Urine detection commonly reflects asymptomatic shedding and does not diagnose PML. Results must be interpreted by the treating specialist with the immune status, symptoms, and other investigations.
\section*{Regulatory Status and Authorization Date}
Regulatory status applies to the specific assay, instrument, manufacturer, and intended use rather than to the test name alone. No single approval date can be assigned to this test category. For a commercial assay, verify the current product in the relevant regulator's database: the U.S. Food and Drug Administration (FDA) clearance, approval, or authorization; India's Central Drugs Standard Control Organisation (CDSCO) licence or permission; the European Union In Vitro Diagnostic Regulation (EU IVDR) CE certificate issued through the applicable conformity-assessment route; or the United Kingdom Medicines and Healthcare products Regulatory Agency (MHRA) registration and UKCA/CE status. Laboratory-developed tests may not have an individual FDA, CDSCO, EU IVDR, or MHRA approval date.
\clearpage
